\documentclass[12pt,a4paper]{article}
\usepackage[utf8]{vietnam}
\usepackage{amsmath,amssymb}
\usepackage{tabularx}
\usepackage{tikz}
\usetikzlibrary{calc,angles,quotes,arrows.meta,patterns}
\usepackage{xcolor}
\usepackage{enumitem}
\usepackage{fancyhdr}
\usepackage{graphicx}
\usepackage[top=1.0cm, bottom=1.0cm, left=1.2cm, right=1.2cm]{geometry}
\usepackage[dethi]{ex_test}
\SAOPTN{widthSA=3.2,heightSA=0.62}

\setlength{\headheight}{15pt}
\setlength{\parskip}{0pt}

\pagestyle{fancy}
\fancyhf{}
\rhead{\textit{Môn Toán 8 -- Đề số 04}}
\lhead{\textit{Bộ đề ôn tập Giữa học kì I}}
\cfoot{\thepage}

\begin{document}

% --- KHUNG TIÊU ĐỀ ĐỀ THI ---
\noindent
\begin{tabularx}{\linewidth}{X r}
\textbf{TRƯỜNG THCS: \dotfill} & \textbf{ĐỀ KIỂM TRA ĐỊNH KỲ GIỮA HỌC KỲ I} \\
\textbf{Họ và tên: \dotfill} & \textbf{NĂM HỌC 2026 -- 2027} \\
\textbf{Lớp: \dotfill \quad SBD: \dotfill} & \textbf{Môn: TOÁN 8} \\
\hline
\end{tabularx}
\vspace{2mm}

\begin{center}
  \textbf{\large ĐỀ ÔN TẬP GIỮA HỌC KỲ I -- ĐỀ SỐ 04}\\
  \textit{Thời gian làm bài: 90 phút (Không kể thời gian phát đề)}
\end{center}
\vspace{1mm}

% --- PHẦN I ---
\noindent\textbf{\large PHẦN I. CÂU TRẮC NGHIỆM NHIỀU PHƯƠNG ÁN LỰA CHỌN (3,0 điểm)}

\textit{Thí sinh trả lời từ câu 1 đến câu 12. Mỗi câu hỏi thí sinh chỉ chọn một phương án duy nhất.}

\vspace{1mm}
\begin{ex}
Trong các biểu thức sau, biểu thức nào là đơn thức?
\choice
{\True $3x^2y$}
{$2x + y$}
{$x^2 - 1$}
{$\dfrac{x}{y}$}
\end{ex}

\begin{ex}
Đơn thức nào sau đây đồng dạng với đơn thức $-5x^2y^3$?
\choice
{$2x^3y^2$}
{\True $7x^2y^3$}
{$-5xy^3$}
{$3x^2y^2$}
\end{ex}

\begin{ex}
Bậc của đa thức $P = 4x^2y^3 - 5x^3y^3 + 2xy - 1$ là:
\choice
{5}
{\True 6}
{4}
{2}
\end{ex}

\begin{ex}
Kết quả của phép nhân $2x \cdot (x^2 - 3y)$ là:
\choice
{$2x^3 - 3y$}
{\True $2x^3 - 6xy$}
{$2x^2 - 6xy$}
{$2x^3 + 6xy$}
\end{ex}

\begin{ex}
Kết quả của phép chia $12x^4y^3 : 3x^2y$ là:
\choice
{\True $4x^2y^2$}
{$4x^2y^3$}
{$4x^6y^4$}
{$9x^2y^2$}
\end{ex}

\begin{ex}
Khai triển hằng đẳng thức $(x - 3y)^2$ ta được kết quả là:
\choice
{\True $x^2 - 6xy + 9y^2$}
{$x^2 - 3xy + 9y^2$}
{$x^2 - 9y^2$}
{$x^2 + 6xy + 9y^2$}
\end{ex}

\begin{ex}
Biểu thức $x^2 - 16$ được viết dưới dạng tích là:
\choice
{$(x - 16)(x + 16)$}
{$(x - 4)^2$}
{\True $(x - 4)(x + 4)$}
{$(x + 4)^2$}
\end{ex}

\begin{ex}
Tổng số đo các góc trong một tứ giác bằng:
\choice
{$180^\circ$}
{\True $360^\circ$}
{$270^\circ$}
{$90^\circ$}
\end{ex}

\begin{ex}
Cho tứ giác $ABCD$ có $\widehat{A} = 80^\circ$, $\widehat{B} = 100^\circ$, $\widehat{C} = 110^\circ$. Số đo góc $D$ là:
\choice
{\True $70^\circ$}
{$80^\circ$}
{$90^\circ$}
{$100^\circ$}
\end{ex}

\begin{ex}
Hình bình hành có một góc vuông là:
\choice
{Hình thoi}
{Hình thang cân}
{\True Hình chữ nhật}
{Hình vuông}
\end{ex}

\begin{ex}
Các mặt bên của hình chóp tam giác đều là:
\choice
{Các tam giác đều}
{\True Các tam giác cân}
{Các tam giác vuông}
{Các hình chữ nhật}
\end{ex}

\begin{ex}
Cho hình chóp tứ giác đều có độ dài cạnh đáy bằng $5\text{ cm}$ và độ dài trung đoạn bằng $6\text{ cm}$. Diện tích xung quanh của hình chóp đó là:
\choice
{$30\text{ cm}^2$}
{\True $60\text{ cm}^2$}
{$120\text{ cm}^2$}
{$15\text{ cm}^2$}
\end{ex}


\vspace{1mm}
\noindent\textbf{BẢNG TRẢ LỜI TRẮC NGHIỆM PHẦN I:}
\begin{center}
\begin{tabular}{|c|c|c|c|c|c|c|c|c|c|c|c|c|}
\hline
\textbf{Câu} & \textbf{1} & \textbf{2} & \textbf{3} & \textbf{4} & \textbf{5} & \textbf{6} & \textbf{7} & \textbf{8} & \textbf{9} & \textbf{10} & \textbf{11} & \textbf{12} \\
\hline
\textbf{Chọn} & \rule{0pt}{0.65cm} & & & & & & & & & & & \\
\hline
\end{tabular}
\end{center}
\vspace{2mm}

\setcounter{ex}{0}
\noindent\textbf{\large PHẦN II. CÂU TRẮC NGHIỆM ĐÚNG SAI (2,0 điểm)}

\textit{Thí sinh trả lời câu 1 và câu 2. Trong mỗi ý a), b), c), d) ở mỗi câu, thí sinh chọn đúng hoặc sai.}

\vspace{1mm}
\begin{ex}
Cho biểu thức $A = (x + 2)^2 + (x - 1)(x + 1) - 2x(x + 2)$. Xét tính đúng sai của các mệnh đề sau:
\choiceTF
{\True Khai triển hằng đẳng thức $(x + 2)^2$ ta được $x^2 + 4x + 4$}
{\True Kết quả của tích $(x - 1)(x + 1)$ là $x^2 - 1$}
{\True Kết quả rút gọn của biểu thức $A$ bằng $3$}
{Giá trị của biểu thức $A$ phụ thuộc vào giá trị của biến $x$}
\end{ex}

\begin{ex}
Cho tam giác $ABC$ vuông tại $A$, gọi $M$ là trung điểm của cạnh huyền $BC$. Từ $M$ kẻ $MD \perp AB$ tại $D$ và $ME \perp AC$ tại $E$. Xét tính đúng sai của các mệnh đề sau:
\choiceTF
{\True Tứ giác $ADME$ có 3 góc vuông nên là hình chữ nhật}
{\True Điểm $D$ là trung điểm của đoạn thẳng $AB$}
{Hai đường chéo $AM$ và $DE$ vuông góc với nhau tại trung điểm mỗi đường}
{\True Cho $AB = 6\text{ cm}$ và $AC = 8\text{ cm}$, độ dài đoạn thẳng $DE$ bằng $5\text{ cm}$}
\end{ex}


\noindent\textbf{BẢNG TRẢ LỜI PHẦN II:}
\begin{center}
\begin{tabular}{|c|c|c|c|c||c|c|c|c|c|}
\hline
\textbf{Câu 1} & \textbf{a)} & \textbf{b)} & \textbf{c)} & \textbf{d)} & \textbf{Câu 2} & \textbf{a)} & \textbf{b)} & \textbf{c)} & \textbf{d)} \\
\hline
\textbf{Đ/S} & \rule{0pt}{0.6cm} & & & & \textbf{Đ/S} & & & & \\
\hline
\end{tabular}
\end{center}
\vspace{2mm}

\setcounter{ex}{0}
\noindent\textbf{\large PHẦN III. CÂU HỎI TRẢ LỜI NGẮN (2,0 điểm)}

\textit{Thí sinh trả lời câu 1 và câu 2. Điền kết quả vào ô trống tương ứng.}

\vspace{1mm}
\begin{ex}
Tính giá trị của biểu thức $M = x^2 - 6x + 9$ tại $x = 103$. \shortans{10000}
\end{ex}

\begin{ex}
Cho hình bình hành $ABCD$ có $\widehat{A} = 110^\circ$. Tính số đo các góc còn lại của hình bình hành? \shortans{$\widehat{B} = 70^\circ, \widehat{C} = 110^\circ, \widehat{D} = 70^\circ$}
\end{ex}


\setcounter{ex}{0}
\noindent\textbf{\large PHẦN IV. TỰ LUẬN (3,0 điểm)}

\textit{Thí sinh trình bày chi tiết lời giải vào giấy làm bài.}

\vspace{1mm}
\begin{ex} (0,75 điểm)
Thực hiện phép tính:
\begin{enumerate}[label=\textbf{\alph*)}, leftmargin=0.8cm, itemsep=2pt]
\item (0,25 điểm) $(2xy^3 - 4y - 8x) \cdot \left(\dfrac{1}{2}y\right)$
\item (0,25 điểm) $(x - y)(x^2 - 2x + y) - x^3 + x^2y$
\item (0,25 điểm) $(-12x^{13}y^{15} + 6x^{10}y^{14}) : (-3x^{10}y^{14})$
\end{enumerate}
\end{ex}
\vspace{0.5mm}

\begin{ex} (0,75 điểm)
Phân tích các đa thức sau thành nhân tử:
\begin{enumerate}[label=\textbf{\alph*)}, leftmargin=0.8cm, itemsep=2pt]
\item (0,25 điểm) $x^2 - 16$
\item (0,25 điểm) $2x^3 - 8x^2y + 8xy^2$
\item (0,25 điểm) $x^2 - 2xy + y^2 + 4x^2y - 4xy^2$
\end{enumerate}
\end{ex}
\vspace{0.5mm}

\begin{ex} (1,0 điểm)
Cho tam giác $ABC$ ($AB < AC$). Gọi $M, N$ lần lượt là trung điểm của các cạnh $BC, AC$. Trên tia đối của tia $MN$ lấy điểm $D$ sao cho $MN = MD$.
\begin{enumerate}[label=\textbf{\alph*)}, leftmargin=0.8cm, itemsep=2pt]
\item (0,5 điểm) Tứ giác $BNCD$ là hình gì? Vì sao?
\item (0,25 điểm) Gọi $I$ là trung điểm của $BN$. Chứng minh ba điểm $A, I, D$ thẳng hàng.
\item (0,25 điểm) Tìm điều kiện của tam giác $ABC$ để tứ giác $BNCD$ là hình thang cân.
\end{enumerate}
\end{ex}
\vspace{0.5mm}

\begin{ex} (0,5 điểm)
Tính nhanh giá trị biểu thức sau: $A = 8 \cdot (3^2 + 1)(3^4 + 1)(3^8 + 1)(3^{16} + 1)$.
\begin{enumerate}[label=\textbf{\alph*)}, leftmargin=0.8cm, itemsep=2pt]
\item (0,5 điểm) Tính nhanh giá trị của $A$.
\end{enumerate}
\end{ex}
\vspace{0.5mm}


\begin{center}
\rule{4cm}{0.4pt}\\
\textit{--- HẾT ---}
\end{center}
\end{document}
